% Generated from locale/tr/content/proof-theory/propositions-as-types/introduction.tex; do not hand-edit.


\olfileid[tr]{int}{pty}{int}

\olsection{Giriş}

Lambda kalkülüsü basit bir hesaplama modelidir. Değişkenlerden kurulan
terimler üzerinde işler. $N$ ile $M$ terimse $(NM)$ de bir terimdir
(``$N$'nin~$M$'ye uygulanması''). $N$ bir terimse $\lambd[x][N]$ de bir
terimdir. $N$,~$x$ değişkenini içeriyorsa $\lambd[x][N]$ içindeki karşılık
gelen $x$ geçişleri $\lambd[x]$ işleci tarafından bağlanır ve artık serbest
değildir. $(\lambd[x][N]M)$ biçimindeki bir terimin $\Subst{N}{M}{x}$'e,
yani $N$ içindeki her serbest $x$ geçişinin $M$ ile değiştirilmesi sonucuna
$\beta$-büzüldüğü söylenir. $N'$; $N$'nin bir alt teriminin
$\beta$-büzülmesiyle elde ediliyorsa $N$ terimi~$N'$'ye $\beta$-dönüşür ve
$N \redone N'$ yazarız. Lambda kalkülüsündeki bir hesaplama $N \redone
N_1 \redone N_2 \redone \dots$ dizisidir. Bu dizi, hiçbir alt terimi
$\beta$-büzülemeyen bir $N_k$ teriminde bitiyorsa bu terime
\emph{normal} ya da $N$'nin normal biçimi denir. Lambda kalkülüsündeki
hesaplamaları böyle diziler olarak düşünün. Normal biçime ulaşılırsa
hesaplama durur ve normal biçim hesaplamanın sonucudur. Bu basit hesaplama
modeli Turing-tamdır: her kısmi özyineli fonksiyon, lambda kalkülüsündeki
bir terimle hesaplanabilir.

Haskell Curry ile William Howard, doğal tümdengelim ve lambda kalkülüsü
arasında yakın bir benzerliği birbirlerinden bağımsız olarak keşfettiler.
Sezgisel olarak $\lambd[x][N]$ terimi bir fonksiyondur: $x$ argümandır;
$N$,~$x$ verildiğinde değerin nasıl hesaplanacağını söyler.
$(\lambd[x][N]M)$ terimi, $\lambd[x][N]$ fonksiyonunu $M$ argümanına
uygular ve $\beta$-büzülmesi bunun nasıl değerlendirileceğini söyler:
$N$ içindeki $x$'i~$M$ ile değiştirin (ve ortaya çıkan terimi
değerlendirmeyi sürdürün). Şimdi bunları sabit bir tanım ve değer kümesine
sahip fonksiyonlar olarak düşünün. $\lambd[x][N]$'nin tanım kümesi $A$,
değer kümesi~$B$ ise $x$ değişkeninin yalnız~$A$'dan argümanlar aldığını,
$N$'nin de~$B$'den değerler ürettiğini düşünmeliyiz. Dolayısıyla
$(\lambd[x][N]M)$ ancak $\lambd[x][N]$, $A \to B$ fonksiyonu ve $M \in A$
ise anlamlı olmalıdır. Bu bilgiyi lambda kalkülüsüne eklersek \emph{tipli}
lambda kalkülüsünü elde ederiz. Tipli lambda kalkülüsünde her
$(\lambd[x][N]M)$ ifadesi değil, yalnız $\lambd[x][N]$ ile $M$'nin
``tipleri'' birbirine uyan ifadeler yasaldır; yani $\lambd[x][N]$'nin tipi
$A \lif B$, $M$'nin tipi~$A$ olmalıdır. $\lambd[x][N]$'nin tipi ancak
$x$'in tipi $A$ ve $N$'nin tipi~$B$ ise $A \to B$'dir. Genel olarak her
$(M_1M_2)$ ifadesi yasal değildir. Yalnız $M_1$'in tipi $A \to B$ ve
$M_2$'nin tipi~$A$ olan $(M_1M_2)$ terimleri yasaldır; bu durumda
$(M_1M_2)$'nin tipi~$B$'dir.

Bu tipleme kuralları artık doğal tümdengelimdeki !!{derivation}s{}e açıkça
karşılık gelir: tipler !!{formula}s ile, !!{derivation}s ise lambda
terimleriyle temsil edilir. Örneğin $!A \lif !B$'nin bir !!a{derivation}{}i,
fonksiyon tipi $!A \lif !B$ olan $\lambd[\typeof{x}{!A}][N]$ terimine
karşılık gelir. Doğal tümdengelimin çıkarım kuralları tipli lambda
kalkülüsündeki tipleme kurallarına karşılık gelir; örneğin:
\begin{prooftree}
  \AxiomC{$\Discharge{!A}{x}$}
  \DeduceC{$!B$}
  \DischargeRule{\Intro\lif}{x}
  \UnaryInfC{$!A \lif !B$}
  \DisplayProof\bottomAlignProof
  \quad\text{şuna karşılık gelir}\quad
  \Axiom$x: !A \fCenter N: !B$
  \UnaryInf$ \fCenter \lambd[\typeof{x}{!A}][N] : !A \lif !B$
\end{prooftree}
Sağdaki kural, $x$'in tipi~$!A$ ve $N$'nin tipi~$!B$ ise
$\lambd[\typeof{x}{!A}][N]$'nin tipinin~$!A \lif !B$ olduğu anlamına gelir.

$\Elim\lif$ kuralı uygulama terimleri için tipleme kuralına karşılık gelir:
\[
  \AxiomC{$!A \lif !B$}
  \AxiomC{$!A$}
  \RightLabel{\Elim\lif}
  \BinaryInfC{$!B$}
  \DisplayProof
  \quad\text{şuna karşılık gelir}\quad
  \Axiom$\fCenter M_1: !A \lif !B$
  \Axiom$\fCenter M_2: !A$
  \BinaryInf$ \fCenter (M_1M_2) : !B$
  \DisplayProof
\]
Bir \Intro\lif{} kuralını hemen bir \Elim\lif{} kuralı izliyorsa
!!{derivation} sadeleştirilebilir:
\begin{gather*}
  \AxiomC{$\Discharge{!A}{x}$}
  \DeduceC{$!B$}
  \DischargeRule{\Intro\lif}{x}
  \UnaryInfC{$!A \lif !B$}
  \AxiomC{}
  \DeduceC{$!A$}
  \RightLabel{\Elim\lif}
  \BinaryInfC{$!B$}
  \DisplayProof
  \quad\redone\quad
  \AxiomC{}
  \DeduceC{$!A$}
  \DeduceC{$!B$}
  \DisplayProof
\end{gather*}
Bu, lambda terimlerinin $\beta$-büzülmesine karşılık gelir:
\[
(\lambd[\typeof{x}{!A}][N]M) \quad\redone\quad
\Subst{N}{M}{x}.
\]

$\land$ kurallarıyla ``çarpım'' tipleri ve $\lor$ kurallarıyla ``toplam''
tipleri arasında da benzer karşılıklar vardır.

Basit tipli lambda kalkülüsündeki terimlerle doğal tümdengelim
!!{derivation}s{}i arasındaki bu karşılığa ``Curry--Howard'', ``program olarak
ispat'' ya da ``tip olarak önerme'' karşılığı denir. !!{formula}s{}in
(önermelerin) tiplere ve ispatların terimlere karşılık gelmesine ek olarak,
karşılıkları şöyle özetleyebiliriz:
\begin{center}
  \begin{tabular}{c c c}
    mantık & program \\
    \hline
    önerme/!!{formula} & tip \\
    ispat & terim \\
    varsayım & değişken \\
    kapatılmış varsayım & bağlı değişken \\
    kapatılmamış varsayım & serbest değişken \\
    $!A \lif !B$ & fonksiyon tipi \\
    $!A \land !B$ & çarpım tipi \\
    $!A \lor !B$ & toplam tipi \\
    $\lfalse$ & boş tip \\
    \hline
  \end{tabular}
\end{center}

Curry--Howard karşılığı, otomatik ispat yardımcıları ile program tip
denetleyicilerinin temel taşlarından biridir; çünkü bir önermeye tanıklık
eden ispatı denetlemek (yukarıda yaptığımız gibi), bir programın (terimin)
bildirilen tipe sahip olup olmadığını denetlemeye denktir.
